%mac
\begin{vidu}
Hai xe lăn $A$ và $B$ có khối lượng lần lượt $m_A = 500\ g$ và $m_B = 250\ g$. Chúng được đặt trên mặt phẳng nằm ngang, tiếp xúc nhau qua một lò xo nén. Ban đầu hệ đứng yên, sau khi đốt dây buộc giữ, lò xo giãn ra rất nhanh và đẩy hai xe rời nhau. Xe $A$ thu được vận tốc $v_A = 1,2\ m/s$ về bên trái. Bỏ qua mọi ma sát và coi thời gian tác dụng lực là rất ngắn. Tính:
\begin{enumerate}
\item Vận tốc của xe $B$ ngay sau khi tách.
\item Độ lớn lực mà lò xo tác dụng lên mỗi xe trong khoảng thời gian tương tác $\Delta t = 0,2\ s$.
\item Gia tốc trung bình của mỗi xe trong khoảng thời gian $\Delta t$.
\end{enumerate}

\loigiai{
\begin{itemize}
\item[1.] Theo định luật III Newton, lực hai xe tác dụng lên nhau có độ lớn bằng nhau nên:
$$ m_A v_A = m_B v_B \Rightarrow v_B = \dfrac{m_A v_A}{m_B} = \dfrac{500 \cdot 1,2}{250} = 2,4\ m/s $$

\item[2.] Tính lực theo định luật II Newton:
$$ F = m_A \cdot a_A = \dfrac{m_A v_A}{\Delta t} = \dfrac{0,5 \cdot 1,2}{0,2} = 3\ N $$
$$ \text{(Vì $F_{AB} = F_{BA}$ nên lực tác dụng lên xe $B$ cũng là } F = 3\ N) $$

\item[3.] Tính gia tốc trung bình của từng xe:
$$ a_A = \dfrac{v_A}{\Delta t} = \dfrac{1,2}{0,2} = 6\ m/s^2 $$
$$ a_B = \dfrac{v_B}{\Delta t} = \dfrac{2,4}{0,2} = 12\ m/s^2 $$
\end{itemize}
}
\end{vidu}

%\loai{Bài toán cơ bản}
\begin{vidu}
Xe lăn $1$ có khối lượng $m_1=400\ g$, có gắn một lò xo. Xe lăn $2$ có khối lượng $m_2$. Ta cho hai xe áp gần vào nhau bằng cách buộc dây để nén lò xo. Khi đốt dây buộc, lò xo dãn ra và sau một thời gian $\Delta t$ rất ngắn, hai xe rời nhau với vận tốc $v_1=1,5\ m / s$; $v_2=1\ m / s$. Bỏ qua ảnh hưởng của ma sát trong thời gian $\Delta t$. Giá trị của $m_2$ là bao nhiêu gam?
\shortans[oly]{600}
%	\choice
%	{\True $m_2=600\ g$}
%	{$m_2=500\ g$}
%	{$m_2=700\ g$}
%	{$m_2=800\ g$}
\haidong{10}
\loigiai{
\begin{itemize}
\item Theo định luật III Niutơn thì sau khi đốt dây, lực tương tác giữa hai vật có độ lớn bằng nhau:
\item 		
\item 		$F_{12}=F_{21}$.
\item 		
\item 		Gọi $m_1$ và $m_2$ là khối lượng các vật,
\item 		
\item 		$a_1$ và $a_2$ là gia tốc tức thời của hai vật. Vận tốc ban đầu (vận tốc trước tương tác) của các vật đều bằng không.
\item 		
\item 		Ta có:
\item 		
\item 		VẬT 1: $F_{21}=m_1 a_1=\dfrac{m_1 v_1}{\Delta t}$
\item 		
\item 		VẬT 2: $F_{12}=m_2 a_2=\dfrac{m 2\ V_2}{\Delta t}$
\item 		
\item 		Từ (1) và (2)
\item 		$\Rightarrow m_1\ V_1=m_2\ V_2$
\item 		
\item 		$\Rightarrow m_2=\dfrac{m_1\ V_1}{~V_2}=\dfrac{400.1,5}{1}=600\ g$
\end{itemize}

}
\haidong{7} \end{vidu}
\begin{vidu}
Cho viên bi $A$ chuyển động tới va chạm vào bi $B$ đang đứng yên, $v_{0A}=4\ m / s$ sau va chạm bi $A$ tiếp tục chuyển động theo phương cũ với $v_A=3\ m / s$, thời gian xảy ra va chạm là $0,4\ s$. Tính gia tốc của hai viên bi, biết $m_{A}=200\ g$, $m_{B}=100\ g$.
\choice
{\True $a_{A}=-2,5\ m / s^2;\  a_{B}=5\ m / s^2$}
{$a_{A}=-3,5\ m / s^2;\ a_{B}=4\ m / s^2$}
{$a_{A}=4,5\ m / s^2;\ a_{B}=6\ m / s^2$}
{$a_{A}=5\ m / s^2;\ a_{B}=3\ m / s^2$}
\haidong{10}
\loigiai{
\begin{itemize}
\item Ta có $a_{A}=\dfrac{v-v_0}{\Delta t}=\dfrac{3-4}{0,4}=-2,5\ m / s^2$

\item Theo định luật III Niu-tơn: $\overrightarrow{F_{AB}}=-\overrightarrow{F_{BA}} \Rightarrow a_{B}=-\dfrac{m_{A} a_{A}}{m_{B}}=-\dfrac{0,2 \cdot(-2,5)}{0,1}=5\left(~m / s^2\right)$

\end{itemize}

}
\haidong{7} \end{vidu}

%\loai{Va chạm vuông góc với tường}
\begin{vidu}%[1P2-4.2G][Loai1]
Quả bóng khối lượng $200\ g$ bay đến đập vào tường theo phương vuông góc với vận tốc $90\ km/h$. Bóng bật trở lại theo phương cũ với vận tốc $54\ km/h$. Thời gian bóng chạm tường là $0,05\ s$. Chọn chiều dương là chiều bóng bật trở lại. Độ lớn của lực do tường tác dụng lên quả bóng là
\choice
{$120\ N$}
{\True $160\ N$}
{$80\ N$}
{$40\ N$}
\haidong{10}\loigiai{
\begin{itemize}
\item Từ đề bài, ta có: $v_1 =  - 90\ km/h = -25\ m/s,\ v_2 = 54\ km/h = 15\ m/s$
\item Gia tốc của quả bóng là: $a = \dfrac{v_2 - v_1}{t} = \dfrac{15 - (-25)}{0,05} = 800\ m/s^2$
\item Lực do tường tác dụng lên quả bóng là: $F = ma = 0,2.800 = 160$ N
\item Vậy của lực do tường tác dụng lên quả bóng có độ lớn bằng 160 N.
\end{itemize}
}
\haidong{7} \end{vidu}
\begin{comment}
%\loai{Hai vật va chạm}
\begin{vidu}%[1P2-4.2G][Loai1]
Hai quả cầu trên mặt phẳng nằm ngang, quả (I) chuyển động với vận tốc 4 m/s đến va chạm với quả cầu (II) đang nằm yên. Sau va chạm hai quả cầu cùng chuyển động theo hướng cũ của quả cầu I với vận tốc 2 m/s. Tỉ số khối lượng $\dfrac{m_1}{m_2}$ của hai quả cầu là
\choice
{\True 1}
{2}
{$\dfrac{1}{4}$}
{$\dfrac{1}{2}$}
\haidong{5}\loigiai{
\begin{itemize}
\item Theo định luật III Niuton, ta có: $m_1\vec {a_1}  =  - m_2\vec {a_2}$ 
\item Gọi $\vec {{v_0}} , \vec v$  là vận tốc quả cầu trước và sau tương tác, $\Delta t$ là thời gian tương tác. 
\item Ta có: $m_1\dfrac{{\vec {v_1}  - \vec {{v_{01}}} }}{{\Delta t}} =  - m_2\dfrac{{\vec {v_2}  - \vec {{v_{02}}} }}{{\Delta t}}$
\item Chọn chiều dương là chiều chuyển động của quả cầu (I)
\item Ta có: $m_1\left( {v_1 - {v_{01}}} \right) =  - m_2\left(v_2 - v_{02}\right)$ $\Rightarrow$ $\dfrac{m_1}{m_2} =  - \dfrac{v_2 - v_{02}}{v_1 - v_{01}} =  - \dfrac{2 - 0}{2 - 4} = 1.$
\end{itemize}
}
\haidong{7} \end{vidu}
\end{comment}
%\loai{Va chạm xiên góc với tường}
\begin{vidu}%[1P2-4.2G][Loai1]
\immini[thm]{Quả bóng khối lượng $200\ g$ bay với vận tốc $72\ km/h$ đến đập vào tường và bật lại với độ lớn vận tốc không đổi. Biết va chạm của bóng với tường tuân theo định luật phản xạ gương (góc phản xạ bằng góc tới) và bóng đến đập vào tường dưới góc tới $30^\circ$, thời gian va chạm là $0,05\ s$. Lực do tường tác dụng lên bóng có độ lớn bằng
\choice
{$125\ N$}
{$146\ N$}
{\True $138\ N$}
{$154\ N$}}{\begin{tikzpicture}[scale=0.6,thick,>=stealth',draw=Mapcolor]
\tkzDefPoints{0/2/m, 0.1/2/n, 0.1/-2/p, 0/-2/q, 0/0/o}
\tkzDefShiftPoint[o](-150:4cm){a}
\tkzDefShiftPoint[o](150:4cm){b}
\tkzDefShiftPoint[o](-150:1.6cm){c}
\tkzDefShiftPoint[o](150:1.6cm){d}
\tkzDefShiftPoint[c](150:1.6cm){e}
\tkzDrawSegments[](m,q)
\fill[pattern=north east lines] (m)--(n)--(p)--(q)--cycle;
\tkzDrawSegments[dashed](o,a o,b)
\draw [Mapcolor,very thick,->](a) -- ($(a)!1.6cm!0:(o)$)node[below]{$\vec{v}$};
\draw [Mapcolor,very thick,->](b) -- ($(b)!1.6cm!180:(o)$)node[above]{$\vec{v}'$};
\shade[ball color=Mapcolor] (a) circle (1.3ex);
\shade[ball color=Mapcolor] (b) circle (1.3ex);
\end{tikzpicture}}
\haidong{10}\loigiai{
\immini{\begin{itemize}
\item Đổi $72\ km/h = 20\ m/s$.
\item Vector gia tốc của quả bóng bàn: $\vec a  = \dfrac{{\vec {v'}  - \vec v }}{{\Delta t}} = \dfrac{{\Delta\vec{v} }}{{\Delta t}}$.
\item Từ hình vẽ $\Rightarrow$ $\Delta v = 2v\cos 30^\circ = 2.20.\cos30^\circ = 20\sqrt{3} \ m/s$. 
\item Lực do tường tác dụng lên bóng có độ lớn là:\\ $\vec F$ = m$\vec a$ = m$\dfrac{{|\Delta\vec{v}|}}{{\Delta t}} = 0,2.\dfrac{{20\sqrt 3 }}{{0,05}} = 138$ N.
\end{itemize}}{\begin{tikzpicture}[scale=1,thick,>=stealth']
\tkzDefPoints{0/2/m, 0.1/2/n, 0.1/-2/p, 0/-2/q, 0/0/o}
\tkzDefShiftPoint[o](-150:4cm){a}
\tkzDefShiftPoint[o](150:4cm){b}
\tkzDefShiftPoint[o](-150:1.6cm){c}
\tkzDefShiftPoint[o](150:1.6cm){d}
\tkzDefShiftPoint[c](150:1.6cm){e}
\tkzDrawSegments[](m,q)
\fill[pattern=north east lines] (m)--(n)--(p)--(q)--cycle;
\tkzDrawSegments[dashed](o,a o,b c,d c,e d,e)
\tkzDrawSegments[blue,very thick,->](o,c o,d)
\tkzDrawSegments[red,very thick,->](o,e)
\tkzDrawPoints[size=3,fill=white](o)
\draw [,blue,very thick,->](a) -- ($(a)!1.6cm!0:(o)$)node[below]{$\vec{v}$};
\draw [,blue,very thick,->](b) -- ($(b)!1.6cm!180:(o)$)node[above]{$\vec{v}'$};
\shade[ball color=blue] (a) circle (1.3ex);
\shade[ball color=blue] (b) circle (1.3ex);
\tkzLabelPoint[above](d){$\vec{v}'$}
\tkzLabelPoint[below,xshift=6pt](c){$-\vec{v}$}
\tkzLabelPoint[left](e){$-\vec{\Delta v}$}
\end{tikzpicture}}
}
\haidong{7} \end{vidu}
